\documentclass[11pt,a4paper]{article}
\usepackage[utf8]{inputenc}
\usepackage[T1]{fontenc}
\usepackage[french]{babel}
\frenchbsetup{StandardLists=true}
\usepackage[left=2cm,right=2cm,top=2cm,bottom=2cm]{geometry}
\usepackage[dvipsnames]{xcolor}
\usepackage{fouriernc}
\usepackage{amsmath,amssymb,amsfonts}
\usepackage{array,tabularx,booktabs,multirow}
\usepackage{colortbl}
\usepackage{graphicx}
\usepackage{mdframed}
\usepackage{enumitem}
\usepackage{fancyhdr}
\usepackage{chemformula}
\usepackage{awesomebox}
\setlength{\aweboxvskip}{0mm}
\usepackage{tcolorbox}
\tcbuselibrary{skins}
\usepackage{fontawesome5}

\definecolor{exoheader}{HTML}{1E5AA0}
\definecolor{exolight}{HTML}{DCE8FF}
\definecolor{warnorange}{RGB}{220,100,0}
\definecolor{problemebg}{HTML}{FFFDE7}

\renewcommand{\arraystretch}{1.8}
\setlength{\extrarowheight}{1mm}

\pagestyle{fancy}
\renewcommand\headrulewidth{1pt}
\fancyhead[L]{Académie de Besançon}
\fancyhead[C]{\textbf{TP --- Synthèse du Slime (réticulation d'un polymère)}}
\fancyhead[R]{Terminale / BTS Chimie}
\fancyfoot[C]{page \thepage}
\fancyfoot[R]{2025 -- 2026}
\fancyfoot[L]{Alexis RUIZ}

\newcommand{\sectitle}[1]{%
  \vspace{6pt}
  \noindent\colorbox{exoheader}{\parbox{\dimexpr\linewidth-2\fboxsep}%
    {\color{white}\bfseries\large\strut\quad #1}}
  \vspace{4pt}
}
\newcommand{\subsectitle}[1]{%
  \vspace{4pt}
  \noindent\colorbox{exolight}{\parbox{\dimexpr\linewidth-2\fboxsep}%
    {\color{exoheader}\bfseries\strut\quad #1}}
  \vspace{2pt}
}
\newcounter{question}
\newcommand{\question}{%
  \stepcounter{question}%
  \noindent\textcolor{exoheader}{\textbf{Question \thequestion.}}\quad
}

\begin{document}

\noindent\textbf{NOM :}\hfill\textbf{Prénom :}\hfill\textbf{Classe :}\hfill\phantom{x}\\[0.2cm]
\hrule\vspace{0.2cm}

\begin{center}
  {\large\bfseries\textcolor{exoheader}{Travaux Pratiques de Chimie des Polymères}}\\[4pt]
  {\LARGE\bfseries Réticulation d'un polymère --- Synthèse du Slime\textsuperscript{\textregistered}}\\[2pt]
  {\small Académie de Besançon --- Belfort}
\end{center}
\hrule\vspace{6pt}

\begin{tcolorbox}[enhanced, colback=problemebg, colframe=exoheader, boxrule=1.5pt,
  title={\faFlask\quad Problématique},
  fonttitle=\bfseries\color{white},
  attach boxed title to top left={yshift=-2mm, xshift=4mm},
  boxed title style={colback=exoheader, colframe=exoheader}]
\textit{Le Slime\textsuperscript{\textregistered} est un matériau aux propriétés mécaniques surprenantes :
il se comporte à la fois comme un \textbf{solide} (il casse si on le brusque) et comme un \textbf{liquide}
(il coule lentement sous son propre poids).
\textbf{Quels sont les mécanismes chimiques qui expliquent ces propriétés inhabituelles ?}}
\end{tcolorbox}

\vspace{6pt}

% ══════════════════════════════════════════════════════════════════════════════
\sectitle{Partie I --- Étude des polymères}
% ══════════════════════════════════════════════════════════════════════════════

\begin{mdframed}[backgroundcolor=exolight, linecolor=exoheader, linewidth=1.2pt]
L'\textbf{alcool polyvinylique} (PVA, $M \approx 100\,000$~g/mol) est un polymère soluble dans l'eau,
constitué de longues chaînes moléculaires (\og spaghettis microscopiques \fg{}).
Le \textbf{tétraborate de sodium} (borax, \ch{B4O7Na2}) est un réticulant : il crée des \textbf{liaisons}
entre les chaînes de PVA, formant un réseau tridimensionnel.
\end{mdframed}

\vspace{4pt}
\question Qu'est-ce qu'un \textbf{polymère} ? Donner la définition et un autre exemple courant. \dotfill

\vspace{0.3cm}\noindent\dotfill

\question Quelle est la différence entre un polymère \textbf{linéaire} et un polymère \textbf{réticulé} ?
Schématiser. \dotfill

\vspace{2cm}

\question Expliquer pourquoi une solution de PVA (sans borax) se comporte comme un liquide visqueux. \dotfill

\vspace{0.3cm}\noindent\dotfill

% ══════════════════════════════════════════════════════════════════════════════
\sectitle{Partie II --- Protocole expérimental}
% ══════════════════════════════════════════════════════════════════════════════

\begin{mdframed}[backgroundcolor=orange!15, linecolor=warnorange, linewidth=2pt]
\textbf{\textcolor{warnorange}{\faExclamationTriangle\quad ATTENTION}} ---
La solution de PVA doit être préparée à \textbf{70--80°C} sous agitation constante.
Ne pas ingérer le Slime obtenu.
\end{mdframed}

\vspace{4pt}
\subsectitle{1. Préparation des solutions}

\begin{enumerate}[leftmargin=1.5cm, label=\textcolor{exoheader}{\textbf{\arabic*.}}]
  \item \textbf{Solution de borax à 4\%} : dissoudre 40~g de borax dans 1~L d'eau (préparée à l'avance).
  \item \textbf{Solution de PVA à 4\%} :
    \begin{itemize}
      \item Peser 4~g de PVA.
      \item Dissoudre dans 100~mL d'eau chaude (70--80°C) sous agitation constante et ininterrompue.
      \item Laisser refroidir.
    \end{itemize}
\end{enumerate}

\vspace{4pt}
\subsectitle{2. Préparation du Slime}

\begin{enumerate}[leftmargin=1.5cm, label=\textcolor{exoheader}{\textbf{\arabic*.}}]
  \item Dans un bécher, verser environ \textbf{10~mL de solution de borax}.
  \item Ajouter éventuellement un colorant ou des paillettes.
  \item Verser \textbf{100~mL de solution de PVA} dans le bécher et mélanger vigoureusement à la baguette.
  \item Sortir le Slime du bécher et le pétrir sur une surface lisse jusqu'à homogénéité.
\end{enumerate}

% ══════════════════════════════════════════════════════════════════════════════
\newpage
\sectitle{Partie III --- Étude des propriétés mécaniques}
% ══════════════════════════════════════════════════════════════════════════════

\subsectitle{1. Observations}

\question Observer et décrire le comportement du Slime dans les situations suivantes :

\begin{center}
\begin{tabularx}{\linewidth}{|>{\bfseries\columncolor{exolight}}m{5cm}|X|}
\hline
\rowcolor{exoheader}
\color{white}\textbf{Situation} & \color{white}\textbf{Observation} \\
\hline
Étirement lent & \\[0.8cm]
\hline
Choc brusque (casser) & \\[0.8cm]
\hline
Suspendu (chute lente) & \\[0.8cm]
\hline
Posé sur une surface & \\[0.8cm]
\hline
Fusion de deux morceaux & \\[0.8cm]
\hline
\end{tabularx}
\end{center}

\vspace{4pt}
\subsectitle{2. Interprétation}

\question Le Slime est-il un solide ou un liquide ? Justifier avec les observations. \dotfill

\vspace{0.3cm}\noindent\dotfill

\question On dit que le Slime est un \textbf{fluide non-newtonien}. Chercher la définition
et expliquer ce terme. \dotfill

\vspace{0.3cm}\noindent\dotfill

\question Comment la quantité de borax ajouté influence-t-elle les propriétés du Slime ?
(Formuler une hypothèse et proposer une expérience pour la tester.) \dotfill

\vspace{0.3cm}\noindent\dotfill

\vspace{0.3cm}\noindent\dotfill

\question Expliquer au niveau moléculaire pourquoi le Slime peut s'étirer sur plusieurs mètres
sans se rompre si l'étirement est suffisamment lent. \dotfill

\vspace{0.3cm}\noindent\dotfill

\vspace{0.3cm}\noindent\dotfill

% ══════════════════════════════════════════════════════════════════════════════
\sectitle{Annexe --- Données}
% ══════════════════════════════════════════════════════════════════════════════

\begin{mdframed}[backgroundcolor=exolight, linecolor=exoheader, linewidth=1.2pt]
\begin{itemize}
  \item \textbf{PVA} (alcool polyvinylique) : $M \approx 100\,000$~g/mol ; polymère soluble à chaud.
  \item \textbf{Borax} (tétraborate de sodium) : \ch{B4O7Na2} ; réticulant ; \textbf{pH} solution à 4\% $\approx$ 9.
  \item Monomère du PVA : alcool vinylique \ch{CH2=CHOH}
  \item Le PVA est obtenu par hydrolyse de l'acétate de polyvinyle.
\end{itemize}
\end{mdframed}

\end{document}
